\section{Methodisches Vorgehen}
\label{sec:methodik}

Dieses Kapitel beschreibt das methodische Vorgehen zur Entwicklung und Untersuchung
des Architekturkonzepts. Ausgehend von der wissenschaftlichen Grundlage werden die
modellbasierte Architekturentwicklung, die methodischen Gestaltungsprinzipien, die
prototypische Konkretisierung und die Evaluationslogik erläutert.

\subsection{Forschungsdesign und Vorgehensmodell}
\label{sec:forschungsdesign}

Die vorliegende Arbeit untersucht einen Ansatz zur kontrollierten Überführung
heterogener Engineering-Bestandsinformationen in SysML-v2-Strukturen.
Der Turing Generator dient als Prototyp, um zentrale Eigenschaften des
entwickelten Konzepts zu operationalisieren und einer systematischen
Untersuchung zugänglich zu machen.

Der Ansatz folgt einer gestaltungsorientierten Vorgehensweise. Designorientierte
Forschung verbindet die Entwicklung eines Artefakts mit der Untersuchung, in
welchem Umfang dieses ein zuvor beschriebenes Problem adressiert
\cite{Wieringa2014,Peffers2007}. Methodisch folgt die Arbeit einer Abfolge aus
wissenschaftlicher Ausgangsbasis, vorläufiger Lösungshypothese,
modellbasierter Architekturentwicklung, prototypischer Konkretisierung und
Evaluation.

Auf Basis der in Kapitel~2 dargestellten wissenschaftlichen Arbeiten wurde
zunächst eine vorläufige Lösungshypothese für die Verarbeitung heterogener
Engineering-Information abgeleitet. Diese definiert erste Verarbeitungsaufgaben
für die Erfassung heterogener Ausgangsdaten, deren Kategorisierung, die
Qualitäts- und Ontologieaufbereitung, die KI-gestützte Musterableitung, die
fachliche Prüfung und Autorisierung sowie die strukturierte Ausgabe. Sie dient als
Ausgangspunkt für die anschließende modellbasierte Architekturentwicklung.
Die methodische Herleitung dieser Ausgangsbasis wird in
Abschnitt~\ref{sec:literaturmethodik} beschrieben.

Die frühe Lösungshypothese ist in
Abbildung~\ref{fig:initial_processing_hypothesis} zusammengefasst.

\begin{figure}[H]
    \centering
    \includegraphics[
        width=\textwidth,
        height=0.72\textheight,
        keepaspectratio
    ]{Graphiken/3.1_Initiale_Loesungshypothese.pdf}
    \caption[Frühe literaturbasierte Lösungshypothese]{
    Frühe literaturbasierte Lösungshypothese des Transformationsprozesses
    (eigene Darstellung).}
    \label{fig:initial_processing_hypothesis}
\end{figure}

Der in Abbildung~\ref{fig:initial_processing_hypothesis} dargestellte
Transformationsprozess wird anschließend schrittweise in Anforderungen,
Funktionen und logische Architekturelemente überführt. Die
Architekturentwicklung erfolgt dabei modellbasiert und top-down. Das
methodische Vorgehen wird in Abschnitt~\ref{sec:architekturmethodik}
beschrieben.

Die prototypische Konkretisierung macht ausgewählte Architekturannahmen in
einem ausführbaren System praktisch überprüfbar. Abschnitt~\ref{sec:prototypische_untersuchung}
beschreibt die hierfür verwendete Vorgehensweise.

Die Evaluation umfasst technische, prozessuale und menschliche
Prüfgegenstände. Abschnitt~\ref{sec:evaluationslogik} beschreibt die
zugrunde liegende Evaluationslogik; Kapitel~6 konkretisiert das Vorgehen und
die jeweiligen Prüfgegenstände.

Die Forschungsfragen aus Abschnitt~\ref{sec:forschungsfragen} bilden über
alle Schritte hinweg den konstanten Bezugsrahmen der Untersuchung. Das
Forschungsdesign verbindet damit wissenschaftliche Ausgangsbasis,
Lösungshypothese, modellbasierte Architekturentwicklung, prototypische
Konkretisierung und Evaluation zu einem zusammenhängenden Vorgehen.

\subsection{Literaturrecherche und Ableitung des Forschungsbedarfs}
\label{sec:literaturmethodik}

Die strukturierte Literaturrecherche bildete die wissenschaftliche
Ausgangsbasis der Arbeit. Sie diente sowohl der Ableitung des in Kapitel~2
dargestellten Forschungsbedarfs als auch der Identifikation erster
Lösungsansätze für die Architekturentwicklung. Ziel war es, relevante
wissenschaftliche Arbeiten an der Schnittstelle von MBSE, SysML,
KI-gestützter Modellgenerierung, Brownfield-Transformation und semantischer
Informationsverarbeitung zu identifizieren und hinsichtlich ihrer Bedeutung
für die Forschungsfragen auszuwerten.

Die Planung der Literaturrecherche orientierte sich an etablierten
Vorgehensweisen für nachvollziehbare Literaturstudien. Zur Strukturierung des
Suchraums wurden die Forschungsfragen in Suchbegriffscluster überführt. Dabei
wurden insbesondere Modellierungssprachen, KI-Technologien, MBSE und
Architecture as Code sowie Brownfield- und Heterogenitätsprobleme
berücksichtigt. Das PICOC-Schema diente als Hilfsrahmen zur Strukturierung von
Untersuchungsgegenstand, Intervention, Vergleichsperspektive, Zielgrößen und
Anwendungskontext \cite{KitchenhamCharters2007}. Der vollständige Literature
Search Plan mit Suchclustern, Suchstrings und Auswahlkriterien ist in
Anhang~\ref{app:lsp} dokumentiert.

Die initiale Suche erfolgte in Scopus, ScienceDirect, IEEE Xplore und der ACM
Digital Library. Ergänzend wurde Scopus AI explorativ genutzt, um über
natürlichsprachliche Fragestellungen weitere potenziell relevante Arbeiten zu
identifizieren. Diese Treffer wurden mit dem bestehenden Suchbestand
abgeglichen und nach denselben Auswahlkriterien bewertet wie die Ergebnisse
der klassischen Suche.

Die Auswahl erfolgte zunächst anhand von Titel und Abstract. Potenziell
relevante Arbeiten wurden anschließend auf Volltextebene geprüft. Die
Dokumentation von Identifikation, Screening und Auswahl orientierte sich an
den Transparenzprinzipien von PRISMA \cite{Page2021}. Der initial
dokumentierte Suchlauf umfasste 28 Publikationen. Weitere themenspezifische
Quellen wurden im Verlauf der Arbeit bei neu auftretenden Fragestellungen
ergänzt. Die zugehörige Dokumentation ist im Literature Search Report in
Anhang~\ref{app:lsr} enthalten.

Die Literatur wurde im Hinblick auf die für die Forschungsfragen relevanten
Lösungsansätze analysiert. Betrachtet wurden insbesondere die Transformation
heterogener Engineering-Information, KI- beziehungsweise LLM-gestützte
Modellbildung, semantische Harmonisierung, formale Modellvalidierung,
Human-in-the-Loop sowie Modularität und Interoperabilität. Aus der
Gegenüberstellung dieser Ansätze wurden der in Kapitel~2 dargestellte
Forschungsbedarf und die in Abschnitt~\ref{sec:forschungsdesign}
beschriebene frühe Lösungshypothese abgeleitet.

\subsection{Methodik der modellbasierten Architekturentwicklung}
\label{sec:architekturmethodik}

Die Entwicklung des Architekturkonzepts erfolgte modellbasiert. Ziel war es,
die Forschungsfragen schrittweise in Anforderungen, Funktionen und logische
Architekturelemente zu überführen. Als Strukturierungsrahmen diente das
RFLP-Prinzip, wobei für die vorliegende Arbeit insbesondere die Ebenen
Requirements, Functional und Logical relevant sind. RFLP unterstützt dabei
die schrittweise Konkretisierung fachlicher Bedarfe über gefordertes
Systemverhalten bis zu den dafür erforderlichen logischen
Architekturelementen \cite{LiLockettLawson2020}.

Den Ausgangspunkt der Modellierung bilden die Forschungsfragen, daraus
abgeleitete Stakeholder-Bedürfnisse sowie fachliche und technische
Randbedingungen. Daraus werden zunächst Anforderungen formuliert, die das
fachlich und technisch erforderliche Systemverhalten beschreiben. Auf dieser
Grundlage werden Funktionen abgeleitet, welche die zur Erfüllung der
Anforderungen erforderlichen Leistungen konkretisieren. Die logische
Architektur ordnet diese Funktionen anschließend fachlichen
Verantwortungsbereichen beziehungsweise logischen Elementen zu.

Die Beziehungen zwischen den Modellierungsebenen werden über die
Architekturentwicklung hinweg erhalten, sodass Herkunft und Weiterführung
einer Anforderung über Funktionen bis zu den zugehörigen Architekturelementen
nachvollziehbar bleiben. Diese Nachverfolgbarkeit zwischen Anforderungen und
daraus abgeleiteten Engineering-Artefakten ist ein wesentliches Prinzip des
Requirements Engineering \cite{ISO29148}. Die explizite Beschreibung von
Architekturelementen und ihren Beziehungen unterstützt zugleich eine
konsistente Darstellung unterschiedlicher Architekturbelange
\cite{ISO42010_2022}.

Die Modellierungsebenen werden während der Architekturentwicklung
wechselseitig auf Konsistenz geprüft. Erkenntnisse auf funktionaler oder
logischer Ebene können dadurch zur Präzisierung zugrunde liegender
Anforderungen zurückgeführt werden. Requirements, Functional und Logical
Architecture werden auf diese Weise schrittweise miteinander abgestimmt.

Für die formale Repräsentation des Architekturmodells wird SysML~v2
verwendet. Die Sprache ermöglicht die maschinenverarbeitbare Beschreibung von
Anforderungen, Verhalten, Strukturen und ihren Beziehungen
\cite{SysMLv2}.

Die modellbasierte Architekturentwicklung konzentriert sich auf die Ebenen
Requirements, Functional und Logical. Damit werden die fachlichen
Verantwortungs- und Informationsverarbeitungsstrukturen beschrieben, die für
die Beantwortung der Forschungsfragen relevant sind.

\subsection{Methodische Gestaltungsprinzipien der KI-gestützten Informationsverarbeitung}
\label{sec:ki_gestaltungsprinzipien}

Für die Architekturentwicklung werden mehrere methodische Gestaltungsprinzipien
aus den in Kapitel~2 betrachteten Eigenschaften KI-gestützter
Informationsverarbeitung abgeleitet. Ausgangspunkt ist die Unterscheidung
zwischen Transformationsschritten, die eine semantische Interpretation
erfordern, und solchen, deren Verhalten durch explizite Regeln beschrieben
werden kann. Large Language Models können sprachlich und strukturell
heterogene Informationen verarbeiten, ihre Ergebnisse benötigen jedoch eine
zusätzliche fachliche Bewertung \cite{Souza.2026,Topcu2025}. Für semantisch
offene Aufgaben werden deshalb probabilistische Verfahren eingesetzt, während
formal beschreibbare Verarbeitungsschritte regelgebunden und reproduzierbar
ausgeführt werden. Regelgebundene Verarbeitung erhöht dabei die
Reproduzierbarkeit und technische Prüfbarkeit; die fachliche Bewertung des
Verarbeitungsergebnisses bleibt ein eigenständiger Prüfschritt.

Ein zweites Prinzip ist die Bindung von Interpretationsergebnissen an die
zugrunde liegenden Informationen. Herkunft, Kontext und relevante
Verarbeitungsschritte müssen erhalten bleiben, damit Ergebnisse überprüft und
gegebenenfalls korrigiert werden können. Provenienzmodelle beschreiben diese
Beziehungen zwischen Informationen, Verarbeitungsaktivitäten und ihrer
Entstehung \cite{MoreauMissier2013,Simmhan2005}. Provenienz wird deshalb als
Bestandteil des Verarbeitungsprozesses berücksichtigt.

Für semantisch offene Informationen werden mehrere bewusst unterschiedliche
Interpretationsperspektiven auf dieselbe quellengebundene Informationsbasis
angewendet. Hierzu werden rollenbasierte Personas genutzt, um unterschiedliche
fachliche Perspektiven abzubilden. Untersuchungen mit mehreren
Modellinstanzen beziehungsweise Agenten zeigen, dass unterschiedliche
Perspektiven zu voneinander abweichenden Interpretations- und
Argumentationsergebnissen führen können \cite{Du2024,TriemDing2024}.
Übereinstimmung und Varianz zwischen den Perspektiven werden daher als
zusätzliche Bewertungsinformation für die nachgelagerte fachliche Prüfung
herangezogen. Die Bewertung erfolgt auf Basis der einzelnen Perspektiven und
ihrer Unterschiede statt über eine reine Mehrheitsentscheidung
\cite{Sunkaraneni2026}.

Ein weiteres Gestaltungsprinzip ist die gezielte Einbindung menschlicher
Entscheidungen. Human-in-the-Loop-Ansätze unterscheiden sich darin, an welchen
Stellen menschliche Kontrolle, Bewertung oder Entscheidung in einen
KI-gestützten Prozess eingebunden wird
\cite{MosqueiraRey2023,Lazaros2026}. Für die vorliegende Arbeit werden
menschliche Entscheidungen insbesondere für die fachliche Bewertung und
Autorisierung interpretations- oder modellierungsbedürftiger Informationen
vorgesehen.

Semantische und ontologische Referenzrahmen ergänzen diese Vorgehensweise,
indem sie Begriffe, Relationen und Strukturzusammenhänge explizit beschreibbar
machen und damit die Harmonisierung und Wiederverwendung heterogener
Engineering-Information unterstützen
\cite{ISO21838_2_2021,OrellanaMandrick2019,Lu2022,Kulvatunyou2022}. Sie
strukturieren den verfügbaren semantischen Bezugsrahmen, während die fachliche
Bewertung der daraus abgeleiteten Informationen separat erfolgt.

Für die weitere Architekturentwicklung ergeben sich daraus vier
übergeordnete Leitlinien: die aufgabengerechte Verbindung probabilistischer
und regelgebundener Verarbeitung, die Bindung von Interpretationsergebnissen
an ihre Informationsgrundlage, die Nutzung mehrerer Perspektiven als
zusätzliche Bewertungsinformation sowie die gezielte Einbindung menschlicher
Entscheidungen. Semantische und ontologische Referenzrahmen unterstützen diese
Prinzipien. Kapitel~4 konkretisiert sie in der entwickelten
Informationsverarbeitungsarchitektur.

\subsection{Prototypische Konkretisierung und Untersuchung des Architekturkonzepts}
\label{sec:prototypische_untersuchung}

Das entwickelte Architekturkonzept wird prototypisch konkretisiert und in
einen ausführbaren Systemprototyp überführt. Ziel ist es, zentrale Eigenschaften
des Konzepts praktisch überprüfbar und einer systematischen Untersuchung
zugänglich zu machen \cite{Wieringa2014,Peffers2007}.

Die prototypische Umsetzung erfolgt auf Grundlage des zuvor entwickelten
Architekturmodells. Erkenntnisse aus Realisierung und technischer Prüfung
werden genutzt, um Implementierungsfehler zu korrigieren und erforderliche
Präzisierungen des Architekturmodells nachvollziehbar vorzunehmen. Der
Prototyp bildet damit die technische Grundlage für die anschließende
Untersuchung des Architekturkonzepts.

Die Implementierung wird in wesentlichen Teilen durch ein Large Language
Model unterstützt. Die KI-Unterstützung wird insbesondere zur Erstellung
und Überarbeitung von Programmcode eingesetzt. Fachliche Spezifikation,
Architekturentscheidungen sowie Bewertung und Freigabe der erzeugten
Änderungen verbleiben beim Autor. Die weitergehende Einordnung dieses
KI-gestützten Entwicklungsprozesses erfolgt in Kapitel~8.

Die technische Realisierung des Prototyps wird in Kapitel~5 beschrieben.
Die zur Untersuchung verwendeten Verifikations- und Evaluationsverfahren
werden in Kapitel~6 dargestellt.

\subsection{Evaluationslogik und Prüfgegenstände}
\label{sec:evaluationslogik}

Die Evaluation verbindet mehrere Prüfebenen, um unterschiedliche Eigenschaften
des entwickelten Architekturkonzepts und seiner prototypischen Realisierung zu
untersuchen. Unterschieden werden technische Verifikation, formative
End-to-End- und Mehrquellenprüfung, eine extern begleitete Live-Demonstration,
Nachverfolgbarkeitsprüfung, Validierung der SysML-v2-Artefakte sowie menschliche
Modellprüfung und Freigabe \cite{Wieringa2014,ISO42030}.

Die technische Verifikation prüft zunächst die spezifizierten und
deterministisch überprüfbaren Eigenschaften der prototypischen Realisierung.
Sie bildet damit die Grundlage für die anschließende Untersuchung des
zusammenhängenden Verarbeitungspfads.

Darauf aufbauend werden die End-to-End-Verarbeitung und die Verarbeitung
mehrerer Quellen formativ untersucht. Ziel ist es zu prüfen, ob das
Architekturkonzept im Zusammenspiel seiner vorgesehenen Verarbeitungsschritte
und Kontrollmechanismen durchgängig operationalisiert werden kann.
Beobachtete Abweichungen werden dokumentiert und hinsichtlich ihrer Ursache
bewertet.

Der konsolidierte Prototypstand wird anschließend in einer extern begleiteten
Live-Demonstration über die reguläre Anwendungsschnittstelle ausgeführt. Die
Bedienung sowie die erforderlichen fachlichen und modellbezogenen
Entscheidungen erfolgen durch den Autor. Die Demonstration ermöglicht damit
die zusammenhängende Beobachtung eines vollständigen Verarbeitungslaufs unter
den vorgesehenen Anwendungsbedingungen.

An die Live-Demonstration schließt sich eine Nachverfolgbarkeitsprüfung an.
Dabei wird anhand der erzeugten und gespeicherten Evidenz rekonstruiert, ob die
verarbeiteten Eingangsinformationen bis zum resultierenden
SysML-v2-Artefakt nachvollzogen werden können. Die Prüfung bezieht sich auf
den tatsächlich durchgeführten Verarbeitungslauf.

Die erzeugten SysML-v2-Artefakte werden zusätzlich innerhalb der verwendeten
Werkzeugkette validiert \cite{SysMLv2,Cibrian.2025b}. Diese
Validierung bewertet insbesondere ihre formale Verarbeitbarkeit. Die
fachliche Bewertung des erzeugten Modellzustands erfolgt anschließend durch
eine menschliche Prüfung mit separater Freigabeentscheidung
\cite{Topcu2025}.

Die einzelnen Prüfebenen liefern damit unterschiedliche Formen von Evidenz
für die Untersuchung des Architekturkonzepts und seiner prototypischen
Realisierung. Kapitel~6 beschreibt die konkreten Prüfgegenstände,
Testverfahren und Akzeptanzkriterien. Die dabei beobachteten Ergebnisse
werden getrennt davon in Kapitel~7 dargestellt.
